\section{Introduction}\label{sec:introduction}

The critical planar Ising model has a conformally invariant interface limit.
For a random ferromagnet, the corresponding question asks whether weak
microscopic disorder disappears from the macroscopic law when the environment
is held fixed. In two dimensions the disorder is marginal, so a scale-by-scale
decay estimate is needed in place of a uniform contraction.
We extend the weak symmetric two-valued result of~\cite{R} to general
bounded iid bond laws. The law need not be preserved by planar duality,
so the proof must identify the critical temperature without assuming
distributional self-duality.

\subsection{Model and statement}
Let $\rho$ be a probability measure supported on $[-1,1]$ such that
\begin{equation}\label{eq:law-assumptions}
 \int x\,\rho(\dd x)=0,
 \qquad \sigma_\rho^2:=\int x^2\,\rho(\dd x)>0.
\end{equation}
On the unoriented nearest-neighbor edges of $\Z^2$, let $(\xi_e)_e$ be
independent with law $\rho$, and set
\begin{equation}\label{eq:bond-model}
 J_e=1+\varepsilon\xi_e,\qquad 0<\varepsilon<1.
\end{equation}
Write $\omega=(\xi_e)_e$ for the environment and $\Prob_\rho$ for its product
law. For a finite spin domain $\Lambda$, inverse temperature $\beta\geq0$,
and prescribed exterior spins $\tau$, the quenched Ising law has weight
\[
 \exp\!\left(\beta\sum_{\{x,y\}:\{x,y\}\cap\Lambda\ne\varnothing}
       J_{\{x,y\}}\sigma_x\sigma_y\right),
 \qquad \sigma_x=\tau_x\quad(x\notin\Lambda),
\]
where the sum runs over nearest-neighbor edges with at least one free endpoint.
Expectations for this law are denoted by $\langle\,\cdot\,\rangle^\tau_{
\Lambda,\beta,J}$. Put $B_m=[-m,m]^2\cap\Z^2$ and define
\begin{equation}\label{eq:critical-definition}
 \beta_c^\rho(\varepsilon)
 =\inf\left\{\beta\geq0:
 \E_\rho\!\left[\lim_{m\to\infty}
       \langle\sigma_0\rangle^+_{B_m,\beta,J}\right]>0\right\}.
\end{equation}
The plus-state limit exists by ferromagnetic monotonicity. Throughout the
paper, averaging a quenched law means integrating it against $\Prob_\rho$;
the environment is never reweighted by a partition function.
We write $K_0=\tfrac12\log(1+\sqrt2)$ for the pure critical coupling.

Let $D\subset\C$ be a bounded Jordan domain with distinct marked boundary
points $a,b$. A lattice approximation is a simply connected square-lattice domain
containing every interior nearest-neighbor bond, embedded with mesh
$\delta$, with a Jordan boundary and marked interface endpoints
$a_\delta,b_\delta$. We call such a domain full. The boundary approximation used
here means that there are homeomorphisms
\begin{equation}\label{eq:boundary-approximation}
 h_n:\partial D\longrightarrow\partial D_n,
 \qquad \sup_{z\in\partial D}|h_n(z)-z|\longrightarrow0,
 \qquad (a_n,b_n)\longrightarrow(a,b).
\end{equation}
The graphs and bond variables live on the unit lattice and are then embedded
by multiplication by $\delta_n$. This convention allows all $n$ to use one
environment. The exterior-neighbor and pinned-boundary-vertex conventions are specified
in Section~\ref{sec:interface}.

Dobrushin conditions fix plus spins on one marked arc and minus spins on
the other. The arc assignments must give a consistent spin at every
pinned site. The spin interface is oriented from $a_n$ to $b_n$ and uses the
deterministic north-east/south-west pairing whenever a contour vertex has
degree four. We use continuous polygonal representatives of the resulting
oriented curves. For two such curves, let
\begin{equation}\label{eq:curve-metric}
 d_{\rm curv}(\gamma,\eta)
   =\inf_{\phi,\psi}\sup_{t\in[0,1]}
          |\gamma(\phi(t))-\eta(\psi(t))|,
 \qquad \bar d_{\rm curv}=1\wedge d_{\rm curv},
\end{equation}
where $\phi,\psi$ range over increasing homeomorphisms of $[0,1]$; curves
at zero distance are identified. For probability laws on this space, set
\begin{equation}\label{eq:BL-definition}
 d_{\rm BL}(\mu,\nu)
 =\sup\left\{\left|\int F\,\dd\mu-\int F\,\dd\nu\right|:
       \|F\|_\infty\leq1,\ 
       |F(\gamma)-F(\eta)|\leq\bar d_{\rm curv}(\gamma,\eta)\right\}.
\end{equation}
This is a bounded-Lipschitz metric for weak convergence in the oriented
uniform-curve topology. Write $Q_{n,\omega}$ for the quenched Dobrushin
interface law and $\mathsf S_{D;a,b}$ for chordal $\SLE_3(D;a,b)$.

\begin{theorem}\label{thm:main}
For every law $\rho$ satisfying \eqref{eq:law-assumptions}, there exists
$\varepsilon_0(\rho)>0$ with the following property. Fix
$0<\varepsilon<\min\{1,\varepsilon_0(\rho)\}$. For every bounded Jordan
domain $D$, distinct $a,b\in\partial D$, and deterministic square-lattice
approximations satisfying \eqref{eq:boundary-approximation} with
$\delta_n\downarrow0$, the interface law at
$\beta=\beta_c^\rho(\varepsilon)$ satisfies
\begin{equation}\label{eq:main-convergence}
 d_{\rm BL}(Q_{n,\omega},\mathsf S_{D;a,b})
     \xrightarrow[n\to\infty]{\Prob_\rho}0.
\end{equation}
The disorder threshold depends on $\rho$ and not on $D$, its marked points,
or the approximation sequence. Moreover,
$\beta_c^\rho(\varepsilon)=\tfrac12\log(1+\sqrt2)+O_\rho(\varepsilon^2)$.
\end{theorem}

\subsection{History and relation to earlier work}
The effect of weak bond disorder is especially delicate in two dimensions.
Harris's relevance criterion places the pure Ising model at the marginal
case, where the specific-heat exponent is zero and the criterion alone
does not decide the effect of randomness~\cite{Harris}.
Dotsenko and Dotsenko's fermionic analysis predicted a double logarithmic
specific-heat singularity~\cite{DD,DDReview}. Renormalization-group work
of Jug, Shalaev, Shankar, and Ludwig developed and refined the description
of logarithmic corrections to pure-model critical behavior
\cite{Jug,Shalaev,Shankar,Ludwig}. The predictions for different correlation
observables were not identical throughout this development. The relevant
physical picture here is marginal irrelevance: weak disorder decreases
under changes of scale, but only logarithmically. These works concern
thermodynamic quantities and correlations; control of a quenched law for
an entire interface also requires geometric estimates.

For the homogeneous critical model, the conformal description is
rigorous. Schramm's stochastic Loewner evolution supplies the candidate
interface laws~\cite{Schramm}. Smirnov established conformal convergence
of the FK-Ising fermionic observable~\cite{Smirnov}, and Chelkak and
Smirnov proved fermionic-observable universality for critical isoradial
Ising models~\cite{ChelkakSmirnov}. Chelkak, Duminil-Copin, Hongler,
Kemppainen, and Smirnov proved that the critical square-lattice spin
interface with Dobrushin boundary conditions converges to chordal
$\mathrm{SLE}_3$ in uniform distance modulo increasing
reparametrization~\cite[Theorem~1 and Equation~(1)]{CDHKS}.
This is the pure interface theorem underlying our comparison argument.

Rigorous universality results for deterministic interactions provide a
second point of comparison. Giuliani, Greenblatt, and Mastropietro proved
universality of infinite-plane energy-correlation scaling limits for
weak finite-range pair perturbations with lattice symmetries
\cite[Theorems~1.1--1.2]{GGM}. For non-integrable spin interfaces,
Greenblatt and Peltola constructed a Grassmann representation of a
discrete martingale and formulated the correlation estimates needed for
a proposed route to $\mathrm{SLE}_3$
\cite[Theorem~4.4 and Conjecture~4.5]{GP}. Their paper does not prove
interface convergence. For random bonds, Mahfouf proved quenched
FK-interface convergence to $\mathrm{SLE}_{16/3}$ in a regime where the
disorder strength tends to zero with the mesh
\cite[Theorem~1.6]{Mahfouf}. Av\'erous and Mahfouf established
high-probability crossing estimates for near-critical random-bond
FK models, also with vanishing disorder for $1<q\le4$
\cite[Theorems~1.5--1.6]{AM}. In the present theorem the nonzero disorder
strength is fixed as the mesh tends to zero, and the interface separates
spin clusters.

The main theorem of~\cite{R} treats symmetric two-valued bond disorder,
whereas its local scale transformation, moment estimates, and marginal
variance estimate are formulated for more general random local potentials.
We verify their hypotheses for the law~$\rho$ and its exact dual, and use
the pure and deterministic boundary-comparison results of~\cite{S}
through their formulations in~\cite{R}. The critical-point argument
adapts the torus method of~\cite[Section~8]{R} to match two separately
tuned trajectories within the same physical temperature family.
This matching replaces distributional self-duality.

Duality has a substantial history in the study of disordered critical
points. Chayes and Shtengel established vanishing magnetization and
critical correlation properties for self-dual disordered ferromagnets
under their stated hypotheses~\cite[Theorems~$1'$--$2'$]{ChayesShtengel}.
Our general one-edge law need not be self-dual. We use Kramers--Wannier
duality~\cite{KW} in the finite-torus form of Bugrij and Shadura, which
relates the four periodic/antiperiodic partition functions for arbitrary
positive inhomogeneous couplings~\cite[Equations~(4)--(5)]{BS}.
The torus comparison identifies the two tuned temperatures; a separate
argument identifies their common value with the spontaneous-magnetization
threshold.

The remaining argument combines classical probabilistic tools with
the geometric inputs of~\cite{R,S}. The Fortuin--Kasteleyn representation~\cite{FK}, the
Edwards--Sokal coupling~\cite{ES}, and the FKG lattice
inequality~\cite{FKG} relate spin barriers, random-cluster connections,
and boundary monotonicity. The decision-tree inequality of O'Donnell,
Saks, Schramm, and Servedio~\cite{OSSS}, extended to monotonic measures
by Duminil-Copin, Raoufi, and Tassion~\cite[Theorem~1.1]{DCRT}, is
applied at each fixed environment before averaging over the bonds.
The geometric use of crossing estimates to control random curves has
precedents in Aizenman--Burchard~\cite{AB} and
Kemppainen--Smirnov~\cite{KS}; Sheffield--Sun~\cite{SS} distinguish
convergence of driving functions from convergence in the uniform curve
topology. The specific passage from boundary barriers and trace
convergence to oriented-curve convergence used here is the progress
criterion of~\cite[Section~9]{R}.


\subsection{The extension and the proof}
The local scale transformation of~\cite{R} rewrites the model as a
sequence of random local perturbations of pure critical Ising.
Section~\ref{sec:scale} states the analytic estimates, and
Section~\ref{sec:tuning} verifies their hypotheses for the primal and
exact dual families and chooses a temperature for each at which the
perturbations tend to zero in environmental moments. The variance coordinate retains local cubic
terms, so the estimates also cover laws with nonzero third moment.

The two families are parametrized by the same physical inverse
temperature, but their selected temperatures need not initially agree.
Section~\ref{sec:matching} compares the four torus partition functions
along these trajectories. Exact duality shows that, at both selected
temperatures, the same primal FK winding event has probability
asymptotic to its pure critical probability. Monotonicity and a quantitative
lower bound for its temperature derivative then force the temperatures to
agree. This identifies a common primal-dual temperature without
assuming symmetry of the disorder law.

Section~\ref{sec:crossings} transfers the decaying perturbations to local
conditional crossing estimates and joint laws of finitely many crossing
tests on tori. The error
terms are geometrically weighted; decay suffices even though the scale
norms need not be summable. Section~\ref{sec:criticality} uses the
estimates on both lattices to identify the common temperature with
\eqref{eq:critical-definition}. Barriers in annuli with disjoint bond
neighborhoods first rule out percolation there. Above that temperature,
an environmental amplification argument and the decision-tree inequality
give rapid dual-arm decay and positive primal magnetization.

Section~\ref{sec:interface} uses the boundary-comparison results
of~\cite{R,S} to compare the random interface with pure interfaces in
slightly displaced domains. The couplings retain the pure comparison
law conditional on the original environment. This conditional marginal
turns the comparison into a quenched statement. A planar sandwich
identifies the trace, and a separate multi-passage estimate controls
its order of traversal, giving the full metric
in~\eqref{eq:curve-metric}.
